% GENERATED FILE. Edit canonical lessons, then run npm run generate:books.
% canonical-source-hash: fnv1a64:14beef9e5f25334c
% canonical-lessons: UR-C442-tajjub, UR-C442-ajib, UR-C442-dang, UR-C442-ajuba, UR-C442-achambha

\chapter{A Surprise}
\label{ch:ur-a-surprise}
\hlchaptermodality{442}

\begin{chapteropening}
\textbf{By the end of this chapter:} \emph{I can say that something amazed me.}

Complete the last lesson of chapter 442: I can say that something amazed me.
\end{chapteropening}

\section[\texorpdfstring{\ur{تعجب} (tājjub)}{تعجب (tājjub)}]{\ur{تعجب} (tājjub) — astonishment}

\label{lesson:UR-C442-tajjub}



\begin{quote}
Before the new one: say the Urdu for brimming over, then the Urdu for a surge of eagerness.
\end{quote}

\subsection*{You'll want to know: \ur{تعجب}}
\textbf{\ur{تعجب}} — \emph{tājjub} — \textquotedblleft{}astonishment\textquotedblright{}.

\begin{quote}
\ur{اس} \ur{خبر} \ur{پر} \ur{سب} \ur{کو} \ur{تعجب} \ur{ہوا۔}
\end{quote}

\begin{quote}
\emph{is khabar par sab ko tājjub huā.}
\end{quote}

\begin{quote}
Everyone was astonished at this news.
\end{quote}

Masculine.

\begin{cousinweb}
\textbf{\ur{تعجب}} (\emph{tājjub}) is usually traced, through Persian, to Arabic \emph{taʿajjub}, \textquotedblleft{}wondering at something\textquotedblright{}.
\end{cousinweb}

\begin{grammarlens}[title={What you've built}]
One more word for describing what you lived and felt.
\end{grammarlens}

\subsection*{Your turn}
\begin{itemize}
\raggedright
  \item \emph{Say it:} \emph{tājjub}
  \item \emph{Say it:} \emph{tājjub}, once more
  \item \emph{Say it:} \emph{is khabar par sab ko tājjub huā.} — Everyone was astonished at this news.
  \item \emph{Recall:} say the Urdu for brimming over, then the Urdu for a surge of eagerness, then say \emph{tājjub} again
\end{itemize}

\begin{tcolorbox}[breakable,colback=teal!4,colframe=teal!35!black,title={Before you move on}]
What does \ur{تعجب} mean? (\textquotedblleft{}Astonishment\textquotedblright{}.) Say it once more.
\end{tcolorbox}
\section[\texorpdfstring{\ur{عجیب} (ajīb)}{عجیب (ajīb)}]{\ur{عجیب} (ajīb) — strange; odd}

\label{lesson:UR-C442-ajib}



\begin{quote}
Before the new one: say the Urdu for a surge of eagerness, then the Urdu for astonishment.
\end{quote}

\subsection*{You'll want to know: \ur{عجیب}}
\textbf{\ur{عجیب}} — \emph{ajīb} — \textquotedblleft{}strange; odd\textquotedblright{}.

\begin{quote}
\ur{یہ} \ur{بات} \ur{بڑی} \ur{عجیب} \ur{ہے۔}
\end{quote}

\begin{quote}
\emph{yih bāt baṛī ajīb hai.}
\end{quote}

\begin{quote}
This is a very strange thing.
\end{quote}

\begin{cousinweb}
\textbf{\ur{عجیب}} (\emph{ajīb}) is usually traced, through Persian, to Arabic \emph{ʿajīb}, \textquotedblleft{}wonderful, strange\textquotedblright{}.
\end{cousinweb}

\begin{grammarlens}[title={What you've built}]
One more word for describing what you lived and felt.
\end{grammarlens}

\subsection*{Your turn}
\begin{itemize}
\raggedright
  \item \emph{Say it:} \emph{ajīb}
  \item \emph{Say it:} \emph{ajīb}, once more
  \item \emph{Say it:} \emph{yih bāt baṛī ajīb hai.} — This is a very strange thing.
  \item \emph{Recall:} say the Urdu for a surge of eagerness, then the Urdu for astonishment, then say \emph{ajīb} again
\end{itemize}

\begin{tcolorbox}[breakable,colback=teal!4,colframe=teal!35!black,title={Before you move on}]
What does \ur{عجیب} mean? (\textquotedblleft{}Strange; odd\textquotedblright{}.) Say it once more.
\end{tcolorbox}
\section[\texorpdfstring{\ur{دنگ} (dang)}{دنگ (dang)}]{\ur{دنگ} (dang) — dumbfounded}

\label{lesson:UR-C442-dang}



\begin{quote}
Before the new one: say the Urdu for astonishment, then the Urdu for strange.
\end{quote}

\subsection*{You'll want to know: \ur{دنگ}}
\textbf{\ur{دنگ}} — \emph{dang} — \textquotedblleft{}dumbfounded\textquotedblright{}.

\begin{quote}
\ur{اس} \ur{کا} \ur{کام} \ur{دیکھ} \ur{کر} \ur{میں} \ur{دنگ} \ur{رہ} \ur{گیا۔}
\end{quote}

\begin{quote}
\emph{us kā kām dekh kar maiṅ dang rah gayā.}
\end{quote}

\begin{quote}
I was dumbfounded when I saw his work.
\end{quote}

\begin{cousinweb}
\textbf{\ur{دنگ}} (\emph{dang}) is usually traced to Persian \emph{dang}, \textquotedblleft{}stunned, stupefied\textquotedblright{}.
\end{cousinweb}

\begin{grammarlens}[title={What you've built}]
One more word for describing what you lived and felt.
\end{grammarlens}

\subsection*{Your turn}
\begin{itemize}
\raggedright
  \item \emph{Say it:} \emph{dang}
  \item \emph{Say it:} \emph{dang}, once more
  \item \emph{Say it:} \emph{us kā kām dekh kar maiṅ dang rah gayā.} — I was dumbfounded when I saw his work.
  \item \emph{Recall:} say the Urdu for astonishment, then the Urdu for strange, then say \emph{dang} again
\end{itemize}

\begin{tcolorbox}[breakable,colback=teal!4,colframe=teal!35!black,title={Before you move on}]
What does \ur{دنگ} mean? (\textquotedblleft{}Dumbfounded\textquotedblright{}.) Say it once more.
\end{tcolorbox}
\section[\texorpdfstring{\ur{عجوبہ} (ajūbā)}{عجوبہ (ajūbā)}]{\ur{عجوبہ} (ajūbā) — a wonder; a marvel}

\label{lesson:UR-C442-ajuba}



\begin{quote}
Before the new one: say the Urdu for strange, then the Urdu for dumbfounded.
\end{quote}

\subsection*{You'll want to know: \ur{عجوبہ}}
\textbf{\ur{عجوبہ}} — \emph{ajūbā} — \textquotedblleft{}a wonder; a marvel\textquotedblright{}.

\begin{quote}
\ur{یہ} \ur{پرانا} \ur{محل} \ur{ایک} \ur{عجوبہ} \ur{ہے۔}
\end{quote}

\begin{quote}
\emph{yih purānā mahal ek ajūbā hai.}
\end{quote}

\begin{quote}
This old palace is a marvel.
\end{quote}

Masculine.

\begin{cousinweb}
\textbf{\ur{عجوبہ}} (\emph{ajūbā}) is usually traced, through Persian, to Arabic \emph{uʿjūba}, \textquotedblleft{}something to wonder at\textquotedblright{}.
\end{cousinweb}

\begin{grammarlens}[title={What you've built}]
One more word for describing what you lived and felt.
\end{grammarlens}

\subsection*{Your turn}
\begin{itemize}
\raggedright
  \item \emph{Say it:} \emph{ajūbā}
  \item \emph{Say it:} \emph{ajūbā}, once more
  \item \emph{Say it:} \emph{yih purānā mahal ek ajūbā hai.} — This old palace is a marvel.
  \item \emph{Recall:} say the Urdu for strange, then the Urdu for dumbfounded, then say \emph{ajūbā} again
\end{itemize}

\begin{tcolorbox}[breakable,colback=teal!4,colframe=teal!35!black,title={Before you move on}]
What does \ur{عجوبہ} mean? (\textquotedblleft{}A wonder; a marvel\textquotedblright{}.) Say it once more.
\end{tcolorbox}
\section[\texorpdfstring{\ur{اچنبھا} (achambhā)}{اچنبھا (achambhā)}]{\ur{اچنبھا} (achambhā) — an amazing thing; astonishment}

\label{lesson:UR-C442-achambha}



\begin{quote}
Before the new one: say the Urdu for dumbfounded, then the Urdu for a wonder.
\end{quote}

\subsection*{You'll want to know: \ur{اچنبھا}}
\textbf{\ur{اچنبھا}} — \emph{achambhā} — \textquotedblleft{}an amazing thing; astonishment\textquotedblright{}.

\begin{quote}
\ur{اس} \ur{کے} \ur{لیے} \ur{سمندر} \ur{ایک} \ur{اچنبھا} \ur{تھا۔}
\end{quote}

\begin{quote}
\emph{us ke liye samandar ek achambhā thā.}
\end{quote}

\begin{quote}
For her the sea was an amazing sight.
\end{quote}

Masculine.

\begin{cousinweb}
\textbf{\ur{اچنبھا}} (\emph{achambhā}) is inherited from \emph{ācchambha}, a form CDIAL reconstructs for \textquotedblleft{}stiffness\textquotedblright{} and so \textquotedblleft{}astonishment\textquotedblright{}, the state of being frozen with surprise (CDIAL 1079).
\end{cousinweb}

\begin{grammarlens}[title={What you've built}]
One more word for describing what you lived and felt.
\end{grammarlens}

\subsection*{Your turn}
\begin{itemize}
\raggedright
  \item \emph{Say it:} \emph{achambhā}
  \item \emph{Say it:} \emph{achambhā}, once more
  \item \emph{Say it:} \emph{us ke liye samandar ek achambhā thā.} — For her the sea was an amazing sight.
  \item \emph{Recall:} say the Urdu for dumbfounded, then the Urdu for a wonder, then say \emph{achambhā} again
\end{itemize}

\begin{tcolorbox}[breakable,colback=teal!4,colframe=teal!35!black,title={Before you move on}]
What does \ur{اچنبھا} mean? (\textquotedblleft{}An amazing thing; astonishment\textquotedblright{}.) Say it once more.
\end{tcolorbox}
